%======================================================================
%  Common request and engine performance measurements
%======================================================================

\section{Performance metrics}
\label{sec:metrics}

Latency describes waiting and output pace; throughput describes a
processing rate.  \tabref{tab:inference-metrics} defines the loading,
phase and request measurements used to interpret inference~\cite{vllm2026}.
Their timing boundaries must be specified explicitly.

Comparisons state model, precision, prompt and output lengths, concurrency,
execution-batch limits and scheduling.  More concurrency can improve
aggregate throughput while increasing waiting.  The median describes
typical latency; p95 is the value at or below which 95\% of observations
fall.  Specify whether observations are requests, per-request averages,
or token intervals.  Memory footprint and hardware activity help explain
results; representation changes also require evaluating numerical quality.

\begin{table}[htbp]
\centering
\small
\tabsetup
\begin{tabularx}{\textwidth}{P{2.7cm}P{1.6cm}Y}
\headrow
\headcell{Metric} & \headcell{Unit} & \headcell{Meaning} \\[3pt]
\tabheadrule
\tabstripeRow
Model loading time
  & ms or s
  & Time from starting model preparation to a declared ready state;
    report whether page population, conversion and packing are included. \\[3pt]
\tabrowsep
\tabplainRow
Time to first token (TTFT)
  & ms or s
  & Time from the declared request start to the first output token.
    Describes the wait before generation becomes visible. \\[3pt]
\tabrowsep
\tabstripeRow
Inter-token latency (ITL)
  & ms
  & Time between successive output tokens.  Describes the pace and
    variability of generation. \\[3pt]
\tabrowsep
\tabplainRow
Time per output token (TPOT)
  & ms/token
  & Mean time per output token after the first, within one request.
    Summarises the generation interval. \\[3pt]
\tabrowsep
\tabstripeRow
End-to-end latency
  & ms or s
  & Time from the declared request start to its completion.
    Describes the wait for the full response. \\[3pt]
\tabrowsep
\tabplainRow
Prefill throughput
  & tokens/s
  & Prompt tokens processed per second of prompt evaluation, using a
    direct phase timer. \\[3pt]
\tabrowsep
\tabstripeRow
Output throughput
  & tokens/s
  & Generated tokens per second over a stated interval, for one
    sequence or aggregated across requests. \\[3pt]
\tabbottomrule
\end{tabularx}
\caption{Common inference performance metrics.  Smaller latency and
larger throughput indicate better performance under otherwise
comparable workloads.}
\label{tab:inference-metrics}
\end{table}
\FloatBarrier

For a request producing \(O>1\) tokens, let \(t_1\) and \(t_O\)
be the times of the first and last output tokens.  Then

\begin{equation}
\label{eq:tpot}
  T_{\mathrm{TPOT}} = \frac{t_O-t_1}{O-1}.
\end{equation}
\eqcaption{Mean generation time per output token, excluding the first.}

With one token per emission event, TPOT averages a request's ITL
intervals; streamed chunks instead measure delivery of several tokens~\cite{vllm2026}.
TPOT is undefined for one-token responses.  Its reciprocal gives the
post-first-token rate when measured in seconds/token, or
\(1000/T_{\mathrm{TPOT}}\) for milliseconds/token.  Aggregate throughput
is not the reciprocal of one request's TPOT; serving rate can also be
reported in completed requests/s.

TTFT can include queueing, input preparation, prefill and delivery,
according to the declared client or engine boundary~\cite{agrawal2025evaluating}.
The first output normally follows prefill, without an extra decode
evaluation.  Direct phase timers specify whether scheduling, cache
management, sampling and synchronization are included.  Decode throughput
counts evaluated positions, whereas output throughput counts emitted
tokens; ordinary generation produces its first output from prefill logits.

Report model loading and warm-up separately from steady-state inference,
including readiness and page-population conditions.  Warm-up can involve
page faults, kernel preparation and initial allocation.  Aggregate output
throughput divides emitted tokens by a declared observation duration,
with explicit treatment of warm-up, unfinished requests and interval
endpoints.  Client and engine timestamps require consistent boundaries.
